\exercisegroup

\begin{enumerate}
\def\labelenumi{\arabic{enumi}.}
\tightlist
\item
  设 S 为符号 P、P-、X、X- 的集合， \(x_1, \ldots, x_n\) ，这些字母 \(x_i\) 权重为0，P和P-权重为1，X、X-权重为2。对于L(S)中的每个词A，我们如下定义A的方差：每个字母 \(x_i\) ，并且符号 P、X 的方差为 0；符号 P⁻、X⁻ 的方差为 1；A 的方差是 A 中出现的符号的方差之和 *（在域 F₂ 中）*（换句话说，如果 A 包含偶数个方差为 1 的符号，则为 0，否则为 1）。
\end{enumerate}

一个有符号阶梯型是L(S)中的一个平衡词A（第一章，附录，第3节），它满足以下两个条件之一：（1）A是其中一个字母 \(x_i\) ； (2) 若 A 具有形式 \(f\mathbf{A}_1\mathbf{A}_2 \ldots \mathbf{A}_p\) (第一章，附录，第3号，命题2的推论2)，其中 \(f\) 是其中一个标志 \(\mathsf{P}, \mathsf{P}^{-}, \mathsf{X}, \mathsf{X}^{-}, p\) 等于1或2，且 \(\mathbf{A}_i\) ( \(1 \leqslant i \leqslant p\) )是平衡词，则每一个 \(\mathbf{A}_i\) 是带符号阶梯型；进一步，若 \(f = \mathsf{X}\) ，那么 \(\mathbf{A}_1\) 并且 \(\mathbf{A}_2\) 的方差为0，且若 \(f = \mathsf{X}^{-}\) ，那么 \(\mathbf{A}_1\) 并且 \(\mathbf{A}_2\) 方差为1。

一个带符号的阶梯型被称为协变的，如果它的方差为0；被称为逆变的，如果它的方差为1。如果在带符号的阶梯型A中，分别将P-和X-替换为P和X，我们得到阶梯型A*（§1，练习1）。阶梯型A*在n个项E₁上的每一个实现， \ldots，Eₙ 称为符号阶梯型 A 在 E₁ 上的一个实现， \ldots，Eₙ，并记为 A(E₁， \ldots，Eₙ）。

让 \(\mathbf{E}_1, \ldots, \mathbf{E}_n, \mathbf{E}_1', \ldots, \mathbf{E}_n'\) 是集合，并令 \(f_i\) 成为...的映射 \(\mathbf{E}_i\) 进入。 \(\mathbf{E}_i'\) ( \(1 \leqslant i \leqslant n\) ）。证明：对每个定义在 \(x_1, \ldots, x_n\) 上的带符号阶梯型 S，我们可以关联一个映射 \(\{f_1, \ldots, f_n\}^{S}\) 满足以下性质：

\begin{enumerate}
\def\labelenumi{(\roman{enumi})}
\item
  如果S是协变的（相应地，逆变的），那么 \(\{f_1, \ldots, f_n\}^{S}\) 是……的映射 \(\mathrm{S}(\mathrm{E}_1, \ldots, \mathrm{E}_n)\) 进入。 \(\mathrm{S}(\mathrm{E}_1', \ldots, \mathrm{E}_n')\) （分别为 \(\mathrm{S}(\mathrm{E}_1', \ldots, \mathrm{E}_n')\) 进入。 \(\mathrm{S}(\mathrm{E}_1, \ldots, \mathrm{E}_n)\) ).
\item
  如果 S 是一个字母 \(x_{i}\) ，那么 \(\{f_1, \ldots, f_n\}^{S}\) 是 \(f_{i}\) .
\item
  若S是PT（相应地，P-T），且若 \(g = \{f_1, \ldots, f_n\}^{\mathrm{T}}\) 是F到 \(\mathbf{F}'\) ，那么 \(\{f_1, \ldots, f_n\}^{S}\) 是 \(g\) 到子集集合的扩展（分别是 \(\overline{g}^1\) 到子集的集合。
\item
  如果 S 是 XTU 或 X-TU，如果 \(\{f_{1},\ldots,f_{n}\}^{T}\) 是一个映射 \(g:F\to F'\) 并且如果 \(\{f_{1},\ldots,f_{n}\}^{U}\) 是一个映射 \(h:G\to G'\) ，那么 \(\{f_{1},\ldots,f_{n}\}^{S}\) 是扩展 \(g\times h:F\times G\to F'\times G'\) .
\end{enumerate}

该映射 \(\{f_1, \ldots, f_n\}^{S}\) 被称为的带符号典范扩张 \(f_1, \ldots, f_n\) 对应于带符号的阶梯型S。若S是一个阶梯型（即，若 \(\mathbb{P}^{-}\) 并且 \(\mathbb{X}^{-}\) 不出现在 S 中），则带符号的典范扩张 \(\{f_1, \ldots, f_n\}^{S}\) 等于 \(\langle f_1, \ldots, f_n \rangle^{S}\) .

证明：如果 \(f_{i}\) 是……的映射 \(\mathbf{E}_i\) 进入。 \(\mathbf{E}_i'\) ，并且如果 \(f_{i}'\) 是……的映射 \(\mathbf{E}_i'\) 进入。 \(\mathbf{E}_i''\) ( \(1 \leqslant i \leqslant n\) ），那么对于协变带符号阶梯类型 S，我们有

\[
\left\{f _ {1} ^ {\prime} \circ f _ {1}, \dots , f _ {n} ^ {\prime} \circ f _ {n} \right\} ^ {\mathrm{S}} = \left\{f _ {1} ^ {\prime}, \dots , f _ {n} ^ {\prime} \right\} ^ {\mathrm{S}} \circ \left\{f _ {1}, \dots , f _ {n} \right\} ^ {\mathrm{S}},
\]

而对于逆变带符号阶梯类型，我们有

\[
\left\{f _ {1} ^ {\prime} \circ f _ {1}, \dots , f _ {n} ^ {\prime} \circ f _ {n} \right\} ^ {\mathrm{S}} = \left\{f _ {1}, \dots , f _ {n} \right\} ^ {\mathrm{S}} \circ \left\{f _ {1} ^ {\prime}, \dots , f _ {n} ^ {\prime} \right\} ^ {\mathrm{S}}.
\]

推导出，如果 \(f_i\) 是 \(\mathbf{E}_i\) 到上 \(\mathbf{E}_i'\) 并且如果 \(f_i'\) 是其逆双射（ \(1 \leqslant i \leqslant n\) ），则 \(\{f_1, \ldots, f_n\}^S\) 是一个双射且 \(\{f_1', \ldots, f_n'\}^S\) 其逆。此外，如果 \(S^*\) 是（无符号）阶梯类型，对应于有符号阶梯类型 \(S\) ，那么 \(\{f_1, \ldots, f_n\}^S\) 等于 \(\langle f_1, \ldots, f_n \rangle^{S^*}\) 或至 \(\langle f_1', \ldots, f_n' \rangle^{S^*}\) 视情况而定 \(S\) 是协变的或逆变的。

\begin{enumerate}
\def\labelenumi{\arabic{enumi}.}
\setcounter{enumi}{1}
\tightlist
\item
  让S是一个符号阶梯型，作用于 \(n + m\) 字母（练习1）。令 \(\Sigma\) 是结构的一个种类，具有 \(x_{1}, \ldots, x_{n}\) 作为主基集，并且 \(A_{1}, \ldots, A_{m}\) 作为辅助基集，其典型特征形式为 \(s \in \mathfrak{P}(S(x_{1}, \ldots, x_{n}, A_{1}, \ldots, A_{m}))\) 。证明我们可以如下为这类结构定义一个 \(\sigma\) -态射的概念：给定 n 个集合 \(E_{1}, \ldots, E_{n}\) 赋予一个U结构的 \(\Sigma\)，n个集合 \(E_{1}^{\prime}, \ldots, E_{n}^{\prime}\) 赋予一个结构 \(U^{\prime}\) 的物种 \(\Sigma\) ，以及映射 \(f_{i}: E_{i} \to E_{i}^{\prime} (1 \leqslant i \leqslant n)\) ，那么 \((f_{1}, \ldots, f_{n})\) 被称为 \(\sigma\) -态射，如果映射 \(f_{i}\) 满足以下条件：
\end{enumerate}

\begin{enumerate}
\def\labelenumi{(\roman{enumi})}
\tightlist
\item
  如果 S 是协变阶梯类型，则
\end{enumerate}

\[
\{f _ {1}, \dots , f _ {n}, I _ {1}, \dots , I _ {m} \} ^ {S} \langle U \rangle \subset U ^ {\prime}.
\]

\begin{enumerate}
\def\labelenumi{(\roman{enumi})}
\setcounter{enumi}{1}
\tightlist
\item
  如果 \(S\) 是一个逆变阶梯类型，则
\end{enumerate}

\[
\left\{f _ {1}, \dots , f _ {n}, I _ {1}, \dots , I _ {m} \right\} ^ {S} \langle U ^ {\prime} \rangle \subset U.
\]

证明通过适当地指定方差，我们可以以这种方式获得序结构、代数结构和拓扑结构的态射定义。

\begin{enumerate}
\def\labelenumi{\arabic{enumi}.}
\setcounter{enumi}{2}
\item
  设 A、B、C 为三个赋予同一种类结构的集合 \(\Sigma\) ，让 \(f\) 设A到B的一个满态射，且令 \(g\) 设B到C的一个态射。证明如果 \(g \circ f\) 是A到C上的一个同构，则 \(g\) 并且 \(f\) 是 同构。
\item
  设A、B、C、D为四个具有相同种类结构的集合。 \(\Sigma\) ，让 \(f\) 设f为从A到B的一个态射， \(g\) 从 B 到 C 的一个态射，并且 \(h\) 从 C 到 D 的一个态射。证明如果 \(g \circ f\) 并且 \(h \circ g\) 都是同构，则 \(f\) , \(g\) ，以及 \(h\) 同构（参见第二章，§3，习题9）。
\item
  设 A，B 为两个赋予结构的集合 \(\mathcal{S}\) , \(\mathcal{S}'\) 同一类型的结构 \(\Sigma\) . 设 \(f\) 是A到B的一个态射，并令 \(g\) 是B到A的一个态射。令M（相应地N）为所有 \(x \in A\) （分别） \(y \in B\) ) 使得 \(g(f(x)) = x\) （分别） \(f(g(y)) = y\) ）。假设 \(\mathcal{S}\) （分别） \(\mathcal{S}'\) ）在 M（相应地 N）上诱导出一个物种结构 \(\Sigma\) 。证明 M 和 N 在赋予这些结构后是同构的。
\item
  让 \(\Sigma\) 设其为具有主基集 A 和辅助集 k 的结构物种，其典型结构 (F, H) 具有如下典型刻画
\end{enumerate}

\[
\mathbf {F} \in \mathfrak {P} ((\mathbf {A} \times \mathbf {A}) \times \mathbf {A}) \times \mathfrak {P} ((\mathbf {A} \times \mathbf {A}) \times \mathbf {A}) \times \mathfrak {P} ((k \times \mathbf {A}) \times \mathbf {A}) \quad \text { 且 } \quad \mathbf {H} \in \mathfrak {P} (\mathbf {A})
\]

且其公理（可以更明确地写出）如下：“F 是带有单位元的交换 \(k\) -代数结构，且 H 是该代数的一个不可约理想 \(\mathbf{A}''\) （我们回顾，在 \(k\) -代数 A 中，一个不可约理想是指一个理想 \(\mathfrak{a}\) 使得，对于任意理想 \(\mathfrak{b}, \mathfrak{c}\) ，满足 \(\mathfrak{b} \cap \mathfrak{c} = \mathfrak{a}\) ，要么 \(\mathfrak{b} = \mathfrak{a}\) 要么 \(\mathfrak{c} = \mathfrak{a}\) ）。若A，A'是两个分别赋予物种 \(\Sigma\) ，我们定义 \(\sigma\) 从A到A'的态射为同态 \(f\) 的 \(k\) -代数结构(F, H)、(F', H')的集合，这些结构将单位元映到单位元，并且满足 \(f(H) \subset H'\) 。给出一个族（ \(\mathcal{G}_{\lambda}\) ）的物种结构的例子 \(\Sigma\) 在一个集合 A 上，使得该结构族存在最小上界（在物种结构的有序集合中）。 \(\Sigma\) 关于A)，但使得这个最小上界不是该族(A \(_{\lambda}\) , \(\mathcal{G}_{\lambda}\) ，Id \(_{\lambda}\) ），其中 A \(_{\lambda}\) 是赋予该结构的集合 A \(\mathcal{G}_{\lambda}\) ，而 Id \(_{\lambda}\) 是恒等映射 A → A \(_{\lambda}\) 。 （考虑一个多项式环 A = k{[}T{]}：A 中的不可约理想是极大理想的幂。若 F 是 \(k\) -代数结构，考虑两个结构 (F, p \(_1\) ) 和 (F, p \(_2\) 的物种 \(\Sigma\) ，其中 p \(_1\) , p \(_2\) 是 A 的不同极大理想；证明这两个结构的最小上界是 (F, (0))，但这不是所讨论族的最初结构。为证明这一点，考虑 \(k\) -子代数 B = k + (p \(_1\) ∩ p \(_2\) ) 的 A，其中 p \(_1\) ∩ p \(_2\) 是极大理想，且注入 B → A。）另外给出一个结构种的例子 \(\Sigma'\) ，以及一个族（ \(\mathcal{G}_{\lambda}'\) ）的物种结构的例子 \(\Sigma'\) 在集合 A' 上，使得该结构族的最大下界存在，但不是该族的最终结构。（在 A 的对偶（视为向量空间）上考虑线性紧致 \(k\) -由该 \(k\) 通过转置得到A的-代数结构，并定义结构的种类 \(\Sigma'\) 通过转置。）* 7. 设 \(\Sigma\) 令结构的种，其具有一个主基集 A，以及作为辅助基集的集合 R，其一般结构由两个字母 (V, φ) 组成，并具有典型的刻画。

\[
\mathbf {V} \in \mathfrak {P} (\mathfrak {P} (\mathbf {A})) \quad \text { 且 } \quad \varphi \in \mathfrak {P} (\mathbf {R} \times \mathbf {A}),
\]

且其公理为（i）公理R\{V\} 关于拓扑结构的种类（§1，第4号，例3），（ii）关系“存在 \(a>0\) 使得φ是区间（关于拓扑V）的连续单射映射的图 {[}0，a{]} 转换为A''。

如果 A， \(A'\) 是两个集合，分别配备了结构 (V， \(\varphi\) )，(V'， \(\varphi'\) ) 的物种， \(\Sigma\) ， \(\sigma\) 到 A 的 \(A'\) -态射被定义为映射 \(f\) 的 \(\mathbf{A}\) 进入。 \(\mathbf{A}'\) 这些是连续的（关于拓扑 \(\mathbf{V}, \mathbf{V}'\) ），并且其图形 \(\mathbf{F}\) 是这样的 \(\mathbf{F} \circ \varphi \subset \varphi'\) 。证明这一关于 \(\sigma\) -态射的概念可以通过练习2中的程序来定义，并且存在一个乘积结构在两个集合的乘积上 \(\mathbf{A}_1, \mathbf{A}_2\) 赋予任意物种结构的 \(\Sigma\) 但给出一个例子，其中在第一投影下的直接像 \(\mathbf{pr}_1\) ，乘积结构的 \(\mathbf{A}_1 \times \mathbf{A}_2\) 不是给定的原始结构 \(\mathbf{A}_1\) （取 \(\mathbf{A}_1\) 为同胚于 \(\mathbf{R}\) ). *

\begin{enumerate}
\def\labelenumi{\arabic{enumi}.}
\setcounter{enumi}{7}
\tightlist
\item
  让 \(\Sigma\) 是结构的种，它有一个单一的（主）基集，其通用结构 \((V, a, b)\) 由三个字母组成，具有典型的特征描述
\end{enumerate}

\[
\mathrm{V} \in \mathfrak {P} (\mathfrak {P} (\mathrm{A})) \quad \text { 且 } \quad a \in \mathrm{A} \quad \text { 且 } \quad b \in \mathrm{A}
\]

且其公理为关系

\[
\mathrm{R} \{\mathrm{V} \} \quad \text { 且 } \quad a \neq b,
\]

其中 \(R\{V\}\) 是拓扑结构类的公理（ \(\S\) 1，第4号，例3）。若A， \(A'\) 是两个集合，分别配备了结构（V，a，b），（ \(V'\) , \(a'\) , \(b'\) 的物种 \(\Sigma\) ， \(\sigma\) 到 A 的 \(A'\) 被定义为连续映射 \(f: A \to A'\) （关于拓扑 V， \(V'\) ) 使得 \(f(a) = a'\) 并且 \(f(b) = b'\) 证明通过替换 \(\Sigma\) 为等价的结构物种，我们可以获得这一概念 \(\sigma\) -态射，按照练习2的程序。证明如果A、B是两个配备了结构S的集合， \(S'\) 的物种 \(\Sigma\) ，则在 \(A \times B\) ，并且该结构在 \(pr_{1}\) （分别） \(pr_{2}\) ) 下是 S（分别地 \(S'\) ）。* 但给出一个例子，其中不存在 \(pr_{1}\) 的截面，该截面是 \(\sigma\) 到 \(A \times B\) 的 A-态射（取 A 为连通的，B 为离散的）。*

* 9. 设 \(\Sigma\) 是除环结构的种。证明对于这种结构种，一种 \(\sigma\) -态射的概念可以通过取除环的态射来定义 \(\mathbf{K}\) 到除环 \(\mathbf{K}'\) 为同态 \(f\) 的 \(\mathbf{K}\) 进入。 \(\mathbf{K}'\) 在编号4的例2的意义下，连同该映射 \(f_0\) 的 \(\mathbf{K}\) 进入。 \(\mathbf{K}'\) 用于哪些 \(f_0(0) = 0\) 并且 \(f_0(x) = 1\) 对于所有 \(x \neq 0\) 在 \(\mathbf{K}\) 证明：对于这个态射概念，我们有以下性质：对每个除环 \(\mathbf{K}\) （任意特征）的除环结构 \(\mathbf{K}\) 在集合上诱导出一个除环结构（与 \(\mathbf{F}_2\) 的除环结构同构） \(\{0, 1\}\) ；并且如果 \(\mathbf{R}\) 是等价关系，其等价类为 \(\{0\}\) 并且 \(\mathbf{K}^* = \mathbf{K} - \{0\}\) ，存在一个商结构（同构于 \(\mathbf{F}_2\) 的结构 \(\mathbf{K}\) 通过关系 \(\mathbf{R}\) .

*10. 设 \(\Sigma\) 令 A 为具有完备阿基米德有序域结构物种的集合。对于每个赋予该物种结构的集合 A， \(\Sigma\) ，让 \(\varphi_{\mathbf{A}}\) 是唯一的同构 \(\mathbf{A}\) 到上 \(\mathbf{R}\) 。如果 \(\mathbf{A}, \mathbf{B}\) 是两个配备了物种结构的集合，证明 \(\Sigma\) 的态射 \(\mathbf{A}\) 进入。 \(\mathbf{B}\) 可以被视为映射 \(f\) 的 \(\mathbf{A}\) 进入。 \(\mathbf{B}\) 使得 \(\varphi_{\mathbf{B}}(f(x)) \geqslant \varphi_{\mathbf{A}}(x)\) 对于所有 \(x \in \mathbf{A}\) 对于这一态射概念，证明存在双射态射但不是同构，尽管该物种 \(\Sigma\) 是单值的。*

\begin{enumerate}
\def\labelenumi{\arabic{enumi}.}
\setcounter{enumi}{10}
\tightlist
\item
  让 \(\Sigma\) 是结构的一个种类（在一个理论 \(\mathfrak{G}\) 中），它只有一个基集。设 \(s \in \mathrm{F}(x)\) 是典型的刻画，并且 \(\mathbf{R}\{x, s\}\) 公理 \(\Sigma\) . 设 \(\mathbf{A}(x)\) 表示物种 \(\Sigma\) 在 \(x\) 上的结构之集. 设 \(\sigma\{x, y, s, t\}\) 是一个满足条件的项 \((\mathrm{MO}_{\mathbf{I}}), (\mathrm{MO}_{\mathbf{II}})\) 以及以下条件：
\end{enumerate}

\((MO_{\mathrm{III}}^{\prime})\) 给定两个集合 \(E, E^{\prime}\) （在一个理论中 \(C^{\prime}\) 该理论强于 \(C\) ) 并赋予结构 \(S, S^{\prime}\) 的物种 \(\Sigma\) 分别地，如果 \(f\) 是……的任意同构 \(E\) 到上 \(E^{\prime}\) ，那么 \(f \in \sigma \{E, E^{\prime}, S, S^{\prime}\}\) .

\begin{enumerate}
\def\labelenumi{(\alph{enumi})}
\tightlist
\item
  让 \(I_{x}\) 是 x 到自身的恒等映射。证明关系 \(Q\{x, s, t\}\) :
\end{enumerate}

\[
s \in \mathrm{A} (x) \quad \text { 且 } \quad t \in \mathrm{A} (x) \quad \text { 且 } \quad \mathrm{I} _ {x} \in \sigma \left\{x, x, s, t \right\} \cap \sigma \left\{x, x, t, s \right\}
\]

是 \(s\) 与 \(t\) 之间在 \(\mathbf{A}(x)\) 上的等价关系. 设 \(\mathbf{B}(x)\) 是商集 \(\mathbf{A}(x)/\mathbf{Q}\) ，并且让 \(\varphi_x\) 是……的典范映射 \(\mathbf{A}(x)\) 到上 \(\mathbf{B}(x)\) 。假设关系 \(s' \in \mathbf{B}(x)\) 是可迁移的（参见附录中确保这一性质成立的条件），并令 \(\Theta\) 表示具有典型刻画的结构的物种 \(s' \in \mathfrak{P}(\mathbf{F}(x))\) 以及公理 \(s' \in \mathbf{B}(x)\) .

\begin{enumerate}
\def\labelenumi{(\alph{enumi})}
\setcounter{enumi}{1}
\tightlist
\item
  让 \(\overline{\sigma}\{x, y, s', t'\}\) 是所有映射的集合 \(f \in \mathcal{F}(x, y)\) 满足以下关系： \(s' \in \mathrm{B}(x)\) 并且 \(t' \in \mathrm{B}(y)\) 并且存在 \(s \in \mathrm{A}(x)\) 并且 \(t \in \mathrm{A}(y)\) 使得 \(s' = \varphi_x(s)\) 并且 \(t' = \varphi_y(t)\) 并且 \(f \in \sigma\{x, y, s, t\}\) “。证明，对于结构的物种 \(\Theta\) ，项 \(\overline{\sigma}\) 满足 \((\mathrm{MO}_{\mathbf{I}})\) , \((\mathrm{MO}_{\mathbf{II}})\) ，以及 \((\mathrm{MO}_{\mathbf{III}})\) 并且我们有
\end{enumerate}

\[
\sigma \left\{x, y, s, t \right\} \subset \overline {{\sigma}} \left\{x, y, \varphi_ {x} (s), \varphi_ {y} (t) \right\}.
\]

